\section{Kimi K3.1 Pro、API提供とウェイト公開予定}
Moonshot AIのKimi K3.1 Proが利用可能になったとして、X上で性能と公開方針が話題になった。総2.8兆・稼働1040億パラメータのMoE、100万トークンのコンテキストを備え、Artificial Analysisの指標では55点、順位は4位と報告されている。Terminal-Benchの65.7\%や高速版の速度向上も紹介された。フルウェイトは11月4日までに公開予定という。サイバー関連の訓練データを除去し、能力を意図的に抑えたとの主張もあるが、その効果は独立検証されていない。関連報道には観測窓外のものが含まれる。
\dailyxsource{Dekrovによる投稿}{https://x.com/Dekrov_Dev/status/2114113237184569725}

\section{MiniMax M3.1 Flashのウェイト公開}
MiniMax M3.1 FlashのフルウェイトがHugging Faceで公開された。総4280億・稼働230億パラメータのマルチモーダルMoEで、100万トークンのコンテキストに対応するとされる。Xでは量子化やセルフホストを巡る議論が進んだ。コミュニティライセンスには、収益が2000万ドルを超える場合に事前承認が必要との条件が報告されている。モデルの規模に加え、実運用での必要メモリ、推論速度、ライセンス適用範囲が採用判断の焦点となる。
\dailyxsource{Rohan Bhanotによる投稿}{https://x.com/RohanBhanotAI/status/2114106011019854041}

\section{DeepInferに対するSGLangとvLLMの最適化報告}
DeepSeekが数日前に公開した推論フレームワークDeepInferについて、観測窓内にSGLangとvLLMの性能報告が投稿された。両プロジェクトは、prefillとdecodeの処理でDeepInferに匹敵、あるいは上回る結果を主張している。64基のGB300環境で約7.4\%改善したとの数値も紹介されたが、測定条件や独立再現は本号では確認していない。公開済みフレームワークを起点に、周辺エコシステムの最適化が加速している点が注目される。
\dailyxsource{SGLangによる投稿}{https://x.com/sgl_project/status/2114071115605877043}

\section{OpenAIの訓練環境隔離とBrotli脆弱性}
OpenAIがフロンティアモデルの訓練環境からライブインターネット接続を遮断したことと、レッドチーム活動中にBrotliライブラリの未認証リモートコード実行脆弱性を発見したことが、X上で引き続き議論された。モデルの能力評価に加え、訓練環境のネットワーク境界や依存ライブラリの安全性が論点となっている。Grokは参照した公式一次投稿を観測窓外と注記しており、本号では窓内の二次議論として扱う。修正状況や他社への波及は今後の確認事項である。
\dailyxsource{OpenAIの公式投稿（Grokの注記では観測窓外）}{https://x.com/OpenAI/status/2113766354386084338}

\section{Nemotron 3.5 Bolt E4B、約4GBでの動作を訴求}
NVIDIAがNemotron 3.5 Bolt E4B 4GBを公開したと報告された。Gemma 4を基に追加学習した小型モデルとされ、約4GBのメモリで動くローカルエージェントや音声アシスタント向けの用途が注目されている。Xには応答速度や拒否率に関する利用者の感想も見られたが、実機条件を揃えた比較結果は本号では確認していない。Ollamaなどの実行環境への対応や、オンデバイスでの実用性能が次の焦点となる。
\dailyxsource{CatsGPTによる投稿}{https://x.com/Tyr_Nippon/status/2114109607081865304}

\section{Mistral Large 4のプレビュー継続}
Mistral Large 4のパブリックプレビューが続いていることと、月末にウェイトが公開されるのではないかという期待がX上で言及された。約1兆パラメータ級のモデルとして紹介されているが、一次発表は観測窓外であり、窓内の直接的な公式X投稿はGrokの収集では確認できなかった。ウェイト公開は確定した出来事ではなく、今後の発表を待つ段階にある。公開時期だけでなく、利用条件や独立した性能評価にも関心が集まる。

\section{Meta Muse Glimmerの技術報告への再言及}
Metaの30BオープンウェイトモデルMuse Glimmerに関する技術レポートが、観測窓内のX投稿で改めて取り上げられた。消費者向けGPU上でのエージェント性能が議論されたものの、Grokが記録したのは二次的な言及であり、公式の一次投稿は観測窓外とされる。モデルの設計や性能について、この日の収集だけで新たな確定情報が得られたわけではない。詳細な技術評価は元のレポートの確認が必要となる。
\dailyxsource{Hekateによる二次的な言及}{https://x.com/iCleanAI/status/2114114271491879380}
